\scanpage{56}
% archive sheet number: 55
\margin{Ref à SGA4 I}

\uline{Catégories accessibles}

\deleted{\unsure{Catégorie filtrante}} \inserted{Ensemble \inserted{pré}ordonné} grand devant $\pi$. Conditions $L_\pi,L$

Foncteur $\pi$-accessible, accessible

Objet $\pi$-accessible, objet accessible d'une catégorie.
\margin{$9,1$ à $9.3$}

Catégorie $\pi$-accessible : Satisfait $L_\pi$, et admet une \inserted{petite} sous-catégorie \inserted{\uline{str.}} génératrice \deleted{telle que} \inserted{dont les} objets sont $\pi$-accessibles

\margin{(\unsure{modifié|vérifier} 9.4?) c'est la bonne notion si $E$ stable par petites $\varinjlim$; sinon les bonnes catégories sont les $\operatorname{Ind}_\pi(C)$ avec $C$ petite\ldots}% slanted

accessible : $\exists\ \pi$ tel que $C$ $\pi$-accessible

\margin{9.6}
\uline{Prop} \deleted{\unsure{Soit $f$ un fonct}} $\underline{\operatorname{Hom}}(E,F)_\pi$, $F$ satisfaisant $L_\pi$

a) stable par \inserted{les} $\varinjlim$ qui sont représentables dans $F$

b) Si $F$ $\pi$-accessible (faiblement) alors pour tte catégorie $J$ avec $\operatorname{card}\operatorname{Fl}J\leq\pi$ et les $\varprojlim$ de type $J$ repr. dans $F$, $\underline{\operatorname{Hom}}(\ )_\pi$ stable par $\varprojlim$ de type $J$.

\deleted{\unsure{Cor. 1 Hom}}

\uline{Cor 1} Cas \uline{accessible}

\uline{Cor 2} \deleted{\unsure{$E_\pi$ est}} Si $E$ est (fblt) \deleted{\unsure{$\pi$}} \inserted{$\pi$-}accessible, \deleted{\unsure{$E_\pi$ est}} pour tt $\pi'\geq\pi$, $E_{\pi'}$ est stable par $\varinjlim$ de type $J$, avec $\operatorname{card} J\leq\pi'$. [Les objets accessibles de $E$ stables par \inserted{petites} $\varinjlim$]

\uline{Cor 3} Si $E$ est accessible, tout objet de $E$ est accessible.

Réciproque \deleted{\unsure{à}} partielle : Si :

\margin{9.11}
\uline{Prop} \deleted{\unsure{Soit}} $E$ catégorie satisfaisant la condition $L$ et :
\deleted{a) $E$ admet une petite sous-catégorie str. gén.}

a) $E$ stable par noyaux de doubles flèches, et $\operatorname{Ker}$ est accessible (p.ex. commute aux \inserted{petites} $\varinjlim$ filtrantes)

b) $E$ stable par produits fibrés,

c) [\inserted{$E$ admet une petite sous-catégorie \deleted{\unsure{strictement|génératrice}} génératrice par épim. stricts universels} \deleted{Toute} famille épimorphique stricte est épim.\unsure{phique|plique} stricte universelle] \deleted{ou encore : les familles épimorphiques strictes sont épimorphiques strictes] ; i.e. si les familles épimorphiques sont épimorphiques strictes]}
\margin{($E$ admet une pet. ss-cat. str. génératrice, \uline{et} (p.ex}

Alors $E$ est accessible (id., en présence de a)) tout objet de $E$ est accessible

\deleted{\unsure{Conditions}}

[\uline{Cor} $E$ catégorie stable par petites $\varinjlim$, \deleted{et telle que, tte famille épimorphique stricte est épimorphique stricte universelle}

Conditions équivalentes

a) $E$ est accessible

b) $E$ admet une petite sous-catégorie \deleted{\unsure{str|génératrice}} génératrice \inserted{par épim. str. universels} \deleted{(comme} $E$ \deleted{\unsure{stable}} petites $\varprojlim$) et le foncteur Ker sur la catégorie des doubles flèches de $E$ est accessible.

\[
 \left[\ \underline{\mathrm{NB}}\quad E\text{ accessible}\Longrightarrow
 \left\{\begin{array}{l}
 \underline{\operatorname{Hom}}(C,E)\text{ accessible pour tte petite catégorie $C$}\,]\\
 \text{\deleted{\unsure{$E$|$\exists$}} ($E$ $\pi$-accessible) $E_\pi$ petite sous-catégorie gén. par épim.}\\
 \qquad\text{stricts universels}\,]
 \end{array}\right.\right.
\]

\uline{Ex.} \deleted{\unsure{D|Dans}} Un topos, \deleted{une catégorie opposée à un topos} une catégorie \deleted{abélienne} à $\varinjlim$ filtrantes exactes et à petite famille génératrice \deleted{l'opposée d'un} sont des catégories accessibles. Mais $(\mathrm{Ens})^\circ$ et $(\mathrm{Ab})^\circ$ ne sont pas accessibles, p.ex. \deleted{\unsure{$\exists$}} parce que dans ces catégories \deleted{\unsure{le fait que}} \deleted{il n'y a pas de \inserted{petite} sous-catégorie génératrice par épim. stricts universels} le $\operatorname{ker}$ n'est pas accessible [i.e. le foncteur coker dans Ens, (Ab) ne commute pas aux produits quelconque : une $\varprojlim$ grande devant $\pi$ d'épim n'est pas néc. un épim.]

\uline{Filtration cardinale} d'une catégorie $E$ \deleted{satisfaisant \unsure{$L$}} par des $\operatorname{Fil}^{\pi}(E)$, $\pi\geq\pi_0$

a) Les $\operatorname{Fil}^{\pi}(E)$ équivalents à petite catégorie

b) \deleted{\unsure{Que|Fil}} $E$ satisfait $L_{\pi_0}$, $\operatorname{Fil}^{\pi}$ stable par $\varinjlim$ filtrantes \deleted{grandes \unsure{au}} avec $I$ grand devant $\pi_0$ et $\operatorname{card}I\leq\pi$

\margin{9.12. (Condition c) renforcée)}
c) Si $\pi'\geq\pi\geq\pi_0$, tt élément $X$ de $\operatorname{Fil}^{\pi'}$ est de la forme $\varinjlim_{i\in I}X_i$. $X_i\in\operatorname{Fil}^\pi$, $I$ grand devant $\pi$, $\operatorname{card}I\leq\pi'$
